

%  Hamzeh_Khanpour_top_FCNC_FCCee_supporting_notes 

\documentclass[11pt,a4paper]{report}

\usepackage[a4paper,margin=2.25cm]{geometry}

\usepackage[T1]{fontenc}
\usepackage[utf8]{inputenc}
\usepackage{lmodern}
\usepackage{microtype}
\usepackage{amsmath,amssymb,bm,mathtools}
\usepackage{graphicx}
\usepackage{xcolor}
\usepackage{booktabs,longtable,array,tabularx,multirow}
\usepackage{hyperref}
\usepackage{enumitem}
\usepackage[most]{tcolorbox}
\usepackage{pdfpages}
\usepackage{caption}
\usepackage{fancyhdr}
\usepackage{titlesec}

\graphicspath{{./}{Figs/}{Presentation_Higgs_performance_meeting19marc_2024/}}

\hypersetup{colorlinks=true,linkcolor=blue!50!black,citecolor=blue!50!black,urlcolor=blue!60!black}
\definecolor{aghgreen}{RGB}{0,100,60}
\definecolor{aghred}{RGB}{168,20,55}
\definecolor{aghgray}{RGB}{70,70,70}
\definecolor{softgreen}{RGB}{235,248,240}
\definecolor{softred}{RGB}{255,238,242}
\definecolor{softblue}{RGB}{235,242,255}
\definecolor{softgray}{RGB}{246,246,246}

\setlist[itemize]{leftmargin=1.2em,itemsep=0.25em,topsep=0.25em}
\setlist[enumerate]{leftmargin=1.4em,itemsep=0.25em,topsep=0.25em}
\newtcolorbox{keymessagebox}[1][]{enhanced,breakable,colback=softgreen,colframe=aghgreen,title={Key message},fonttitle=\bfseries,#1}
\newtcolorbox{definitionbox}[1][]{enhanced,breakable,colback=softblue,colframe=blue!55!black,title={Definition},fonttitle=\bfseries,#1}
\newtcolorbox{warningbox}[1][]{enhanced,breakable,colback=softred,colframe=aghred,title={Caution},fonttitle=\bfseries,#1}
\newtcolorbox{talknotebox}[1][]{enhanced,breakable,colback=softgray,colframe=aghgray,title={Talk note},fonttitle=\bfseries,#1}
\newtcolorbox{questionbox}[1][]{enhanced,breakable,colback=white,colframe=orange!80!black,title={Likely question},fonttitle=\bfseries,#1}

\newcommand{\ee}{e^+e^-}
\newcommand{\pp}{pp}
\newcommand{\ttbar}{t\bar t}
\newcommand{\br}{\mathrm{Br}}
\newcommand{\CLs}{\mathrm{CL}_s}
\newcommand{\GeV}{\,\mathrm{GeV}}
\newcommand{\abinv}{\,\mathrm{ab}^{-1}}
\newcommand{\fb}{\,\mathrm{fb}}
\newcommand{\ord}{\mathcal{O}}


\setlength{\headheight}{14pt}

\pagestyle{fancy}


\fancyhf{}

\lhead{Precision Physics at the FCC-ee}

\rhead{Hamzeh Khanpour}

\cfoot{\thepage}

\title{\bfseries Precision Physics at the FCC-ee\\[0.4em] Top FCNC Search at the FCC-ee}

\author{Hamzeh Khanpour \\ \\ (hamzeh.khanpour@agh.edu.pl) \\ \\ Faculty of Physics and Applied Computer Science \\ AGH University of Krakow}

\date{Seminarium HEP Bialasowka; 15 May 2026}


\begin{document}
	
	
\maketitle

%\begin{abstract}
	
%This document is a technical companion to the presentation ``Top FCNC search at the FCC-ee: From precision physics at the FCC-ee to rare top interactions''.  It reviews the scientific logic of the slide deck, expands the physics formulae and coupling conventions, explains the analysis strategy, and provides slide-by-slide speaking notes.  It is intended for private speaker preparation and for technical follow-up discussions with a particle-physics audience.

%\end{abstract}

\tableofcontents
\clearpage



\part{Precision Physics at the FCC-ee}

\chapter{Physics background and Precision Physics at the FCC-ee}


\section{FCC-ee as a precision machine}

The Future Circular Collider programme is designed as an integrated precision-plus-energy-frontier strategy.  The electron--positron stage, FCC-ee, would operate at several centre-of-mass energies chosen to maximise sensitivity to different sectors of the Standard Model: the Z pole, the WW threshold, the $ZH$ Higgs stage, the $\ttbar$ threshold and a high-energy 365 GeV stage.  The hadron stage, FCC-hh, then uses the same infrastructure to extend direct reach and high-energy precision.  The relevant point for this talk is that rare top searches at FCC-ee are not disconnected from the rest of the programme: they exploit the same clean event reconstruction, vertexing, flavour tagging and global interpretation logic that define the Higgs and electroweak programmes \cite{FCCFSR,FCCManifesto,FCCCDR}.

\begin{definitionbox}
	
A precision observable is an experimentally measurable quantity whose Standard-Model prediction and experimental uncertainty can be controlled well enough that small deviations become meaningful.  Examples include $m_Z$, $\Gamma_Z$, $\sin^2\theta_{\rm eff}^\ell$, $\sigma(ZH)$, $m_W$, $m_t$ and rare branching fractions such as $\br(t\to q\gamma)$.

\end{definitionbox}

In an effective-field-theory picture, the sensitivity of a precision observable to heavy new physics can be estimated schematically as
\begin{equation}
  \frac{\delta\mathcal O}{\mathcal O_{\rm SM}}
  \sim C_i\frac{v^2}{\Lambda^2},
  \label{eq:v2scale}
\end{equation}
for observables controlled by electroweak symmetry breaking, or as
\begin{equation}
  \frac{\delta\mathcal O}{\mathcal O_{\rm SM}}
  \sim C_i\frac{E^2}{\Lambda^2},
  \label{eq:e2scale}
\end{equation}
for high-energy amplitudes whose deviations grow with the scattering energy $E$.  Here $C_i$ is a dimensionless Wilson coefficient, $v\simeq246\GeV$ is the Higgs vacuum expectation value and $\Lambda$ is the characteristic new-physics scale.  A per-mille measurement therefore probes multi-TeV scales for order-one Wilson coefficients.  This simple estimate is only a guide; real interpretations require correlations, theory uncertainties, detector effects and EFT validity checks.

\begin{keymessagebox}
The first scientific message of the talk is that precision measurements and rare processes are complementary discovery tools.  Precision Higgs, electroweak and top measurements test correlated deviations, while top-FCNC searches test interactions that are essentially absent in the SM.
\end{keymessagebox}

\section{Why Z, WW, Higgs and top measurements are correlated}
The FCC-ee programme is best understood as a network of measurements.  The Z pole fixes electroweak vertices and pseudo-observables.  The WW threshold tests the charged-current gauge sector and the W mass.  The Higgs stage measures $\sigma(ZH)$ inclusively and exclusive Higgs decay rates.  The top threshold and 365 GeV stage constrain $m_t$, $\Gamma_t$ and top electroweak couplings.  In the Standard Model Effective Field Theory (SMEFT), the same gauge-invariant operators can contribute to several of these sectors simultaneously:
\begin{equation}
  \mathcal L_{\rm SMEFT}=\mathcal L_{\rm SM}+
  \sum_i\frac{C_i^{(6)}}{\Lambda^2}\mathcal O_i^{(6)}+
  \sum_i\frac{C_i^{(8)}}{\Lambda^4}\mathcal O_i^{(8)}+\cdots .
\end{equation}
A single-operator interpretation of one Higgs rate is usually too naive, because electroweak data may constrain the same direction.  Conversely, Higgs rates can help interpret electroweak and top deviations.  This is the reason for Slide 13: Higgs precision needs the Z and WW programmes.

\section{Electroweak pseudo-observables and $Zf\bar f$ couplings}
At the Z pole, the collider scans the resonance line shape and measures partial widths and asymmetries.  The most important pseudo-observables include
\begin{equation}
  m_Z,\quad \Gamma_Z,\quad \sigma^0_{\rm had},\quad R_\ell,\quad R_b,\quad R_c,
  \quad A_{\rm FB}^{0,f},\quad \sin^2\theta_{\rm eff}^\ell,\quad N_\nu.
\end{equation}
The hadronic pole cross section is
\begin{equation}
  \sigma^0_{\rm had}=\frac{12\pi}{m_Z^2}\frac{\Gamma_e\Gamma_{\rm had}}{\Gamma_Z^2}.
\end{equation}
The ratios
\begin{equation}
  R_\ell=\frac{\Gamma_{\rm had}}{\Gamma_{\ell\ell}},\qquad
  R_b=\frac{\Gamma_{b\bar b}}{\Gamma_{\rm had}},\qquad
  R_c=\frac{\Gamma_{c\bar c}}{\Gamma_{\rm had}}
\end{equation}
remove some common systematics and isolate flavour-dependent electroweak couplings.  The effective number of light neutrinos can be expressed schematically as
\begin{equation}
  N_\nu=\frac{\Gamma_{\rm inv}}{\Gamma_{\nu}^{\rm SM}},
\end{equation}
where $\Gamma_{\rm inv}$ is the invisible Z width after subtracting the visible partial widths.  Deviations in $N_\nu$ may arise from new invisible particles or modified neutrino couplings, but precision interpretation requires careful control of the full Z-width fit.

\subsection{Chiral couplings}
The Z boson couples differently to left- and right-handed fermions.  Define the chirality projectors
\begin{equation}
  P_L=\frac{1-\gamma_5}{2},\qquad P_R=\frac{1+\gamma_5}{2},
  \qquad f_L=P_L f,\quad f_R=P_R f.
\end{equation}
A useful convention is
\begin{equation}
  \mathcal L_{Zff}=\frac{g}{c_W}Z_\mu\bar f\gamma^\mu
  \left(g_L^fP_L+g_R^fP_R\right)f,
\end{equation}
with
\begin{equation}
  g_L^f=T_3^f-Q_f\sin^2\theta_{\rm eff}^f,
  \qquad
  g_R^f=-Q_f\sin^2\theta_{\rm eff}^f.
\end{equation}
Right-handed SM fermions are singlets under $SU(2)_L$, so $T_3=0$ for their weak isospin.  Left-handed fermions sit in $SU(2)_L$ doublets and have $T_3=\pm 1/2$.  The vector and axial combinations are conventionally related by
\begin{equation}
  g_V^f=g_L^f+g_R^f,
  \qquad
  g_A^f=g_L^f-g_R^f,
\end{equation}
up to global normalisation conventions.  Partial widths and asymmetries probe different combinations:
\begin{equation}
  \Gamma_f\propto (g_L^f)^2+(g_R^f)^2,
  \qquad
  A_f=\frac{(g_L^f)^2-(g_R^f)^2}{(g_L^f)^2+(g_R^f)^2},
\end{equation}
\begin{equation}
  A_{\rm FB}^{0,f}=\frac{3}{4}A_eA_f.
\end{equation}
This is why rates, line-shape information and asymmetries together separate left- and right-handed shifts.

\section{WW threshold, gauge structure and aTGCs}
Near the WW threshold the cross section $\ee\to W^+W^-$ is very sensitive to $m_W$ and $\Gamma_W$.  Above threshold, differential angular distributions are sensitive to the non-Abelian electroweak gauge structure.  A common anomalous triple-gauge-coupling parameterisation is
\begin{align}
  \mathcal L_{WWV}=&-i g_{WWV}\Big[
  g_1^V\left(W^+_{\mu\nu}W^{-\mu}V^\nu-W^-_{\mu\nu}W^{+\mu}V^\nu\right)
  +\kappa_V W^+_\mu W^-_\nu V^{\mu\nu} \\
  &\hspace{3.0cm}
  +\frac{\lambda_V}{m_W^2}W^+_{\mu\nu}W^{-\nu\rho}V_{\rho}^{\ \mu}\Big],
  \qquad V=\gamma,Z.
\end{align}
At tree level in the SM,
\begin{equation}
  g_1^\gamma=g_1^Z=\kappa_\gamma=\kappa_Z=1,
  \qquad
  \lambda_\gamma=\lambda_Z=0.
\end{equation}
The slide uses $\delta g_{1,Z}$, $\delta\kappa_\gamma$ and $\lambda_Z$ as the key deviations.  Here $\delta g_{1,Z}$ modifies the overall $WWZ$ structure, $\delta\kappa_\gamma$ acts like an anomalous magnetic-dipole contribution to $WW\gamma$, and $\lambda_Z$ is a momentum-dependent higher-derivative interaction.  In SMEFT, these parameters are correlated with Higgs-gauge interactions such as $HZZ$ and $HWW$.

\section{Higgs recoil method and absolute coupling extraction}
At $\sqrt{s}\simeq240\GeV$, the process
\begin{equation}
  \ee\to ZH
\end{equation}
operates near the Higgsstrahlung cross-section maximum.  If the Z boson is reconstructed, for example through $Z\to\mu^+\mu^-$ or $Z\to e^+e^-$, the recoil mass is
\begin{align}
  m_{\rm recoil}^2&=(p_{e^+}+p_{e^-}-p_Z)^2 \\
  &=s+m_Z^2-2\sqrt{s}E_Z,
\end{align}
where $E_Z$ is the reconstructed Z-boson energy.  The key point is that this observable does not require reconstructing the Higgs decay.  Therefore invisible, exotic or untagged Higgs decays still contribute to the inclusive $ZH$ peak.  The inclusive cross section scales as
\begin{equation}
  \sigma(ZH)\propto g_{HZZ}^2\equiv \kappa_Z^2(g_{HZZ}^{\rm SM})^2.
\end{equation}
This gives the absolute Higgs normalisation that is missing at hadron colliders.

Exclusive Higgs decays then measure products of production and decay factors:
\begin{equation}
  \sigma(ZH)\,\br(H\to X)
  \propto \frac{g_{HZZ}^2g_{HXX}^2}{\Gamma_H},
\end{equation}
while WW-fusion gives
\begin{equation}
  \sigma(\nu\bar\nu H)\,\br(H\to X)
  \propto \frac{g_{HWW}^2g_{HXX}^2}{\Gamma_H}.
\end{equation}
Combining recoil, exclusive decays and vector-boson fusion gives direct access to absolute couplings and the total width $\Gamma_H$.

\subsection{Exclusive Higgs decay channels}
\begin{longtable}{p{0.20\textwidth}p{0.24\textwidth}p{0.42\textwidth}}
\toprule
Mode & Coupling information & Physics role \\
\midrule
$H\to b\bar b$ & bottom Yukawa $\kappa_b$ & Dominant Higgs decay; central for absolute coupling fits and flavour-tagging performance. \\
$H\to c\bar c$ & charm Yukawa $\kappa_c$ & Direct second-generation quark coupling; requires charm tagging and control of $b/c/g$ separation. \\
$H\to gg$ & effective $Hgg$ modifier $\kappa_g$ & Loop-induced in the SM; sensitive to top loops, QCD and possible coloured BSM states. \\
$H\to \tau^+\tau^-$ & tau Yukawa $\kappa_\tau$ & Tests third-generation lepton mass generation and tau reconstruction. \\
$H\to WW^*$ & $HWW$ coupling $\kappa_W$ & Important width handle; one W is off shell for a 125 GeV Higgs. \\
$H\to ZZ^*$ & $HZZ$ coupling $\kappa_Z$ & Clean gauge-coupling channel and width handle. \\
$H\to\gamma\gamma$ & effective $H\gamma\gamma$ modifier $\kappa_\gamma$ & Rare loop-induced mode; FCC-hh statistics strongly improve precision. \\
$H\to Z\gamma$ & effective $HZ\gamma$ modifier $\kappa_{Z\gamma}$ & Rare loop-induced mode with strong BSM sensitivity. \\
$H\to\mu^+\mu^-$ & muon Yukawa $\kappa_\mu$ & Second-generation lepton Yukawa; benefits strongly from FCC-hh. \\
\bottomrule
\end{longtable}

\section{Higgs coupling modifiers and the $\kappa$-3 framework}
In the coupling-modifier framework,
\begin{equation}
  \kappa_i=\frac{g_i}{g_i^{\rm SM}},
\end{equation}
for tree-level couplings, while loop-induced processes are assigned effective modifiers such as $\kappa_g$, $\kappa_\gamma$ and $\kappa_{Z\gamma}$.  A production-times-decay rate can be written as
\begin{equation}
  (\sigma\cdot \br)(i\to H\to f)=\sigma_i\frac{\Gamma_f}{\Gamma_H}.
\end{equation}
In the $\kappa$ framework,
\begin{equation}
  (\sigma\cdot \br)(i\to H\to f)
  =\left(\sigma_i^{\rm SM}\kappa_i^2\right)
  \frac{\Gamma_f^{\rm SM}\kappa_f^2}{\Gamma_H^{\rm SM}\kappa_H^2},
\end{equation}
so that the signal strength is
\begin{equation}
  \mu_i^f=\frac{(\sigma\cdot \br)}{(\sigma\cdot \br)_{\rm SM}}
  =\frac{\kappa_i^2\kappa_f^2}{\kappa_H^2}.
\end{equation}
The total-width modifier is approximated by
\begin{equation}
  \kappa_H^2=\sum_j \kappa_j^2\br_j^{\rm SM},
\end{equation}
with the sum over all relevant SM decay modes.  If invisible or untagged modes are allowed, the $\kappa$-3 framework often uses
\begin{equation}
  \bar\kappa_H^2=\frac{\kappa_H^2}{1-B_{\rm inv}-B_{\rm und}},
  \qquad
  \Gamma_H=\bar\kappa_H^2\Gamma_H^{\rm SM}.
\end{equation}
The condition $\kappa_V\leq1$ is sometimes imposed in hadron-collider fits to stabilise the width interpretation.  At FCC-ee, inclusive $ZH$ recoil strongly reduces this degeneracy by directly measuring $\sigma(ZH)$.

\begin{warningbox}
The $\kappa$ framework is useful for explaining Higgs-rate measurements, but it is not a gauge-invariant EFT.  SMEFT is required when one wants to connect Higgs, electroweak, diboson and top observables consistently.
\end{warningbox}

\section{SMEFT interpretation of Higgs, EW and top data}
The SMEFT operators most relevant for the motivation slides include bosonic operators and current operators.  Examples are
\begin{align}
  \mathcal O_{HWB}&=H^\dagger\sigma^aH W^a_{\mu\nu}B^{\mu\nu},\\
  \mathcal O_{HD}&=\left|H^\dagger D_\mu H\right|^2,\\
  \mathcal O_{H\ell}^{(1)}&=(H^\dagger i\overleftrightarrow D_\mu H)(\bar\ell\gamma^\mu\ell),\\
  \mathcal O_{H\ell}^{(3)}&=(H^\dagger i\overleftrightarrow D_\mu^a H)(\bar\ell\sigma^a\gamma^\mu\ell),\\
  \mathcal O_{He}&=(H^\dagger i\overleftrightarrow D_\mu H)(\bar e\gamma^\mu e).
\end{align}
These can shift $Zee$ couplings, affect $HZZ$ and $HZee$ contact interactions and correlate with diboson and Higgs production.  Bosonic operators such as $\mathcal O_{HW}$, $\mathcal O_{HB}$ and $\mathcal O_W$ affect gauge-boson kinetic/momentum structures and anomalous triple-gauge couplings.  This is the origin of the chords in the circos plot: the plot visualises correlations among fitted effective couplings, not independent measurements.

\section{Top-quark precision at FCC-ee}
The top threshold programme measures the final state
\begin{equation}
  \ee\to W^+bW^-\bar b,
\end{equation}
with resonant and non-resonant contributions.  The threshold line shape can be written schematically as
\begin{equation}
  \sigma_{WbWb}=\sigma_{WbWb}(\sqrt{s};m_t,\Gamma_t,\alpha_s,y_t).
\end{equation}
The horizontal position of the line shape is sensitive to the top mass, the width controls the sharpness, $\alpha_s$ affects the QCD potential, and the top Yukawa coupling produces smaller but relevant effects through Higgs exchange.  A threshold mass such as the potential-subtracted mass $m_t^{\rm PS}$ is used to avoid the pole-mass ambiguity.

Above threshold, the process
\begin{equation}
  \ee\to \gamma^*,Z^*\to t\bar t
\end{equation}
probes the top electroweak vertex.  A compact notation is
\begin{equation}
  \mathcal L_{ttZ}=\frac{g}{2c_W}\bar t\gamma^\mu\left(g_V^t-g_A^t\gamma_5\right)tZ_\mu.
\end{equation}
Angular distributions and forward--backward asymmetries are essential to separate chiral structures.

\section{FCNC in the Standard Model and beyond}
Flavour-changing neutral currents (FCNCs) such as $t\to q\gamma$, $t\to qZ$, $t\to qH$ and $t\to qg$ are absent at tree level in the SM.  They appear only through loops and are strongly suppressed by the Glashow--Iliopoulos--Maiani mechanism.  The expected branching ratios are tiny, often around $10^{-12}$ or below depending on the channel.  Therefore any observable rate far above the SM expectation would strongly indicate BSM dynamics \cite{AguilarSaavedra2004}.

The talk focuses on $tq\gamma$ and $tqZ$ because these interactions induce clean single-top production in $\ee$ collisions:
\begin{equation}
  \ee\to \gamma^*/Z^*\to t\bar q,\ \bar t q,
  \qquad q=u,c.
\end{equation}
The produced top decays through the SM mode
\begin{equation}
  t\to Wb,
\end{equation}
which leads to semileptonic and fully hadronic signatures.

\section{Effective Lagrangian for top-FCNC interactions}
The anomalous-coupling framework used in the slides follows the common Aguilar-Saavedra normalisation \cite{AguilarSaavedra2009}.  The interaction is
\begin{align}
\mathcal L_{\rm FCNC}^{tqX}=\sum_{q=u,c}\Bigg[&
\frac{g_s}{2m_t}\bar q\lambda^a\sigma^{\mu\nu}(\zeta^L_{qt}P_L+\zeta^R_{qt}P_R)tG^a_{\mu\nu}
-\frac{1}{\sqrt2}\bar q(\eta^L_{qt}P_L+\eta^R_{qt}P_R)tH \\
&-\frac{g_W}{2c_W}\bar q\gamma^\mu(X^L_{qt}P_L+X^R_{qt}P_R)tZ_\mu
+\frac{g_W}{4c_Wm_Z}\bar q\sigma^{\mu\nu}(\kappa^L_{qt}P_L+\kappa^R_{qt}P_R)tZ_{\mu\nu}\\
&+\frac{e}{2m_t}\bar q\sigma^{\mu\nu}(\lambda^L_{qt}P_L+\lambda^R_{qt}P_R)tA_{\mu\nu}
\Bigg]+{\rm h.c.}
\label{eq:fcncLag}
\end{align}
Here $q=u,c$, $g_s$ is the strong coupling, $g_W$ the weak $SU(2)_L$ coupling, $e$ the electromagnetic coupling, $c_W=\cos\theta_W$, $\lambda^a$ are colour generators and $\sigma^{\mu\nu}=i[\gamma^\mu,\gamma^\nu]/2$.  The field strengths are $G^a_{\mu\nu}$, $Z_{\mu\nu}$ and $A_{\mu\nu}$.  The coefficients $\zeta$, $\eta$, $X$, $\kappa$ and $\lambda$ parameterise, respectively, $tqg$, $tqH$, vector $tqZ$, dipole $tqZ$ and dipole $tq\gamma$ interactions.

\begin{warningbox}
The benchmark assumption $\lambda^L=\lambda^R$ or $\kappa^L=\kappa^R$ is a practical signal-generation choice.  It is not a model-independent statement.  Different BSM models can generate strongly chiral patterns.
\end{warningbox}

\section{Translation between anomalous couplings and branching ratios}
The branching ratio is defined as
\begin{equation}
  \br(t\to qX)=\frac{\Gamma(t\to qX)}{\Gamma_t^{\rm SM}+\sum_X\Gamma(t\to qX)}
  \simeq\frac{\Gamma(t\to qX)}{\Gamma_t^{\rm SM}},
\end{equation}
where the approximation holds for small FCNC widths.  In the convention of Eq.~\eqref{eq:fcncLag}, neglecting the light-quark mass, the photon-dipole partial width is commonly written as
\begin{equation}
  \Gamma(t\to q\gamma)=\frac{\alpha}{4}m_t\left(|\lambda^L_{qt}|^2+|\lambda^R_{qt}|^2\right).
\end{equation}
For the Z dipole, with $r_Z=m_Z^2/m_t^2$,
\begin{equation}
  \Gamma(t\to qZ)_\kappa=\frac{\alpha}{32s_W^2c_W^2}\frac{m_t^3}{m_Z^2}(1-r_Z)^2(2+r_Z)
  \left(|\kappa^L_{qt}|^2+|\kappa^R_{qt}|^2\right).
\end{equation}
For the symmetric benchmark $\lambda^L=\lambda^R=\lambda$ and $\kappa^L=\kappa^R=\kappa$,
\begin{align}
  \br(t\to q\gamma)&\simeq\frac{\alpha m_t}{2\Gamma_t^{\rm SM}}|\lambda_{qt}|^2,\\
  \br(t\to qZ)_\kappa&\simeq\frac{\alpha}{16s_W^2c_W^2}\frac{m_t^3}{m_Z^2\Gamma_t^{\rm SM}}(1-r_Z)^2(2+r_Z)|\kappa_{qt}|^2.
\end{align}
The numerical translation must use the same coupling normalisation as the UFO model used for event generation.

\section{Signal processes at 240 and 365 GeV}
The signal process is
\begin{equation}
  \ee\to t\bar q+\bar t q,
\end{equation}
mediated by the anomalous $tq\gamma$ and $tqZ$ couplings.  At 240 GeV, the stage is optimised for Higgs studies but also provides clean single-top kinematics.  At 365 GeV, there is more phase space and harder objects, but also genuine top backgrounds.  The benchmark luminosities used in the slides are $5\abinv$ at 240 GeV and $1.5\abinv$ at 365 GeV.

\section{Semileptonic and fully hadronic final states}
The semileptonic topology is
\begin{equation}
  \ee\to tq\to Wbq\to \ell\nu bq,
\end{equation}
with one isolated electron or muon, missing momentum from the neutrino, one $b$-tagged jet and one energetic light jet.  It is clean but has a smaller branching fraction.  The fully hadronic topology is
\begin{equation}
  \ee\to tq\to Wbq\to q\bar q' bq,
\end{equation}
with four exclusive jets.  It has a larger branching fraction but suffers from larger diboson backgrounds and combinatorics.  A typical assignment variable is
\begin{equation}
  \chi^2=\frac{(m_{jj}-m_W)^2}{\sigma_W^2}+\frac{(m_{bjj}-m_t)^2}{\sigma_t^2},
\end{equation}
where the jet combination minimising $\chi^2$ is chosen as the top candidate.

\section{Backgrounds, jet clustering and flavour tagging}
The dominant backgrounds are $WW$, $ZZ$, $ZH$, $e\gamma$ initiated processes, and at 365 GeV also $\ttbar$ and single-top-like final states.  The analysis adopts the ee-$k_T$ Durham algorithm in exclusive two-jet or four-jet modes.  E-scheme recombination means four-momenta are added directly.  The slides quote a flavour-tagging working point $\epsilon_b=80\%$ and $\epsilon_c=10\%$.  In the FCNC selection, $b$ tagging identifies the top decay, while light-jet rejection and charm/light separation control backgrounds.

Discriminating variables include reconstructed masses $M_W$ and $M_{\rm top}$, the $b$-jet momentum $p_b$, the light-jet momentum $p_{\rm light}$, missing momentum, angular separations and event activity.  If $H_T$ is used, define it in the analysis note as the scalar sum of visible transverse momenta or the corresponding FCC-ee event-activity proxy:
\begin{equation}
  H_T=\sum_i p_{T,i}.
\end{equation}
The usual hadron-collider angular distance is
\begin{equation}
  \Delta R=\sqrt{(\Delta\eta)^2+(\Delta\phi)^2},
\end{equation}
although in $\ee$ analyses ordinary opening angles are often equally transparent.


\section{Statistical analysis and CL$_s$ extraction}
A generic binned likelihood for the combined analysis is
\begin{equation}
  \mathcal L(\mu,\boldsymbol\theta)=\prod_{c,b}{\rm Pois}\left(n_{cb}\,\middle|\,\mu s_{cb}(\boldsymbol\theta)+b_{cb}(\boldsymbol\theta)\right)
  \prod_k\pi_k(\theta_k),
\end{equation}
where $c$ labels channels, $b$ labels histogram bins, $n_{cb}$ are observed or Asimov counts, $s_{cb}$ and $b_{cb}$ are signal and background templates, $\mu$ is the signal-strength parameter, and $\theta_k$ are nuisance parameters.  The CL$_s$ ratio is conventionally
\begin{equation}
  CL_s=\frac{p_{s+b}}{1-p_b},
\end{equation}
with the exclusion set at $CL_s=0.05$ for a 95\% CL upper limit.  In projected studies the observed dataset is often an Asimov dataset, so the quoted limit is an expected sensitivity.




\section{Interpretation of projected limits}
The final tables quote limits on $\br(t\to c\gamma)$ and $\br(t\to cZ)$ rather than directly on $\lambda$ and $\kappa$.  This is useful for comparing with LHC results and with other future-collider projections.  The lepton channel gives the best sensitivity in the present benchmark because the final state is cleaner and the backgrounds are more controllable.  The hadronic channel still contributes because its branching fraction is larger and it probes a complementary topology.  The results should be described as study-level projections, subject to the stated luminosities, detector assumptions, signal model and systematic-treatment assumptions.



\part{Technical appendices}
\appendix
\chapter{Coupling and notation dictionary}
\begin{longtable}{p{0.22\textwidth}p{0.68\textwidth}}
\toprule
Symbol & Meaning \\
\midrule
$\sqrt{s}$ & Centre-of-mass energy of the $\ee$ collision. \\
$m_Z$, $\Gamma_Z$ & Z-boson mass and total width. \\
$m_W$, $\Gamma_W$ & W-boson mass and total width. \\
$\sin^2\theta_{\rm eff}^\ell$ & Effective weak mixing angle measured in leptonic asymmetries. \\
$R_b$ & Ratio $\Gamma_{b\bar b}/\Gamma_{\rm had}$. \\
$N_\nu$ & Effective number of light neutrino species inferred from invisible Z width. \\
$A_{\rm FB}^{0,f}$ & Pole forward--backward asymmetry for fermion $f$. \\
$P_L,P_R$ & Chiral projectors $(1\mp\gamma_5)/2$. \\
$g_L^f,g_R^f$ & Left- and right-handed Z couplings to fermion $f$. \\
$g_V^f,g_A^f$ & Vector and axial-vector Z couplings. \\
$\delta g_{1,Z}$ & Deviation in the $WWZ$ triple-gauge coupling. \\
$\delta\kappa_\gamma$ & Deviation in the $WW\gamma$ magnetic-dipole-like coupling. \\
$\lambda_Z$ & Higher-derivative anomalous $WWZ$ coupling. \\
$\kappa_Z,\kappa_W$ & Higgs coupling modifiers for $HZZ$ and $HWW$. \\
$\kappa_b,\kappa_c,\kappa_t$ & Higgs Yukawa modifiers for bottom, charm and top. \\
$\kappa_g,\kappa_\gamma,\kappa_{Z\gamma}$ & Effective loop-induced Higgs modifiers. \\
$B_{\rm inv}$ & Invisible Higgs branching ratio. \\
$B_{\rm und}$ or $B_{\rm unt}$ & Undetected/untagged Higgs branching ratio. \\
$\Gamma_H$ & Total Higgs width. \\
$m_t^{\rm PS}$ & Potential-subtracted top mass used for threshold fits. \\
$\zeta_{qt}$ & Chromomagnetic $tqg$ FCNC coefficient. \\
$\eta_{qt}$ & Scalar $tqH$ FCNC coefficient. \\
$X_{qt}$ & Vector $tqZ$ FCNC coefficient. \\
$\kappa_{qt}$ & Dipole $tqZ$ FCNC coefficient. \\
$\lambda_{qt}$ & Dipole $tq\gamma$ FCNC coefficient. \\
$\epsilon_b,\epsilon_c$ & $b$-tagging efficiency and charm-tagging/mistag working-point number as used in the analysis. \\
$\CLs$ & Modified frequentist confidence-level ratio used for expected upper limits. \\
\bottomrule
\end{longtable}

\chapter{Electroweak formula sheet}
\begin{align}
\sigma^0_{\rm had}&=\frac{12\pi}{m_Z^2}\frac{\Gamma_e\Gamma_{\rm had}}{\Gamma_Z^2},&
R_\ell&=\frac{\Gamma_{\rm had}}{\Gamma_{\ell\ell}},&
R_b&=\frac{\Gamma_{b\bar b}}{\Gamma_{\rm had}},\\
A_f&=\frac{(g_L^f)^2-(g_R^f)^2}{(g_L^f)^2+(g_R^f)^2},&
A_{\rm FB}^{0,f}&=\frac{3}{4}A_eA_f,&
N_\nu&=\frac{\Gamma_{\rm inv}}{\Gamma_\nu^{\rm SM}}.
\end{align}
These formulae are sufficient for the oral discussion.  A full FCC electroweak fit also includes QED radiator functions, beam-energy calibration, luminosity theory and correlations among pseudo-observables.

\chapter{Higgs recoil and coupling-extraction formula sheet}
\begin{align}
  m_{\rm recoil}^2&=s+m_Z^2-2\sqrt{s}E_Z,\\
  \sigma(ZH)&\propto \kappa_Z^2,\\
  \sigma(ZH)\br(H\to X)&\propto\frac{\kappa_Z^2\kappa_X^2}{\Gamma_H},\\
  \sigma(\nu\bar\nu H)\br(H\to X)&\propto\frac{\kappa_W^2\kappa_X^2}{\Gamma_H},\\
  \mu_i^f&=\frac{\kappa_i^2\kappa_f^2}{\kappa_H^2},\\
  \bar\kappa_H^2&=\frac{\kappa_H^2}{1-B_{\rm inv}-B_{\rm und}}.
\end{align}

\chapter{SMEFT operator map}
\begin{longtable}{p{0.25\textwidth}p{0.55\textwidth}p{0.15\textwidth}}
\toprule
Operator & Main effect \\
\midrule
$\mathcal O_{HWB}$ & Shifts electroweak mixing and correlates EWPOs with Higgs-gauge couplings.  \\
$\mathcal O_{HD}$ & Affects custodial-sensitive electroweak observables and Higgs-gauge interactions.\\
$\mathcal O_{H\ell}^{(1)}$, $\mathcal O_{H\ell}^{(3)}$ & Shift lepton electroweak currents, including $Zee$ and $We\nu$. \\
$\mathcal O_{He}$ & Shifts right-handed electron coupling to the Z.  \\
$\mathcal O_{HW}$, $\mathcal O_{HB}$ & Modify Higgs-gauge and diboson interactions.  \\
$\mathcal O_W$ & Produces momentum-dependent triple-gauge-coupling deviations.  \\
Top dipoles and currents & Affect top EW couplings and loop-induced Higgs/top correlations.  \\
\bottomrule
\end{longtable}

\chapter{Top-FCNC branching-ratio formula sheet}
For the convention used in the slides,
\begin{align}
\Gamma(t\to q\gamma)&=\frac{\alpha}{4}m_t(|\lambda^L_{qt}|^2+|\lambda^R_{qt}|^2),\\
\Gamma(t\to qZ)_\kappa&=\frac{\alpha}{32s_W^2c_W^2}\frac{m_t^3}{m_Z^2}(1-r_Z)^2(2+r_Z)(|\kappa^L_{qt}|^2+|\kappa^R_{qt}|^2),\\
\br(t\to qX)&\simeq\Gamma(t\to qX)/\Gamma_t^{\rm SM},\qquad r_Z=m_Z^2/m_t^2.
\end{align}
\begin{warningbox}
These expressions must not be combined with a differently normalised FCNC Lagrangian.  If another UFO model is used, the conversion between coupling limits and branching ratios must be rederived or checked against the model documentation.
\end{warningbox}

\chapter{Analysis checklist}
\begin{enumerate}
\item Generate FCNC signal with one anomalous coupling switched on at a time.
\item Simulate detector response in the official FCC chain or a documented Delphes/IDEA approximation.
\item Reconstruct exclusive two-jet and four-jet topologies with the selected ee-Durham algorithm.
\item Apply lepton, jet, angular and flavour-tagging selections.
\item Build signal and background templates for discriminating variables.
\item Construct channel-wise likelihoods and combine channels/energies.
\item Translate signal-strength limits into branching-ratio limits using the same Lagrangian normalisation.
\item Repeat with systematic uncertainties and possibly multivariate discriminants for a final public analysis.
\end{enumerate}


\chapter{Extended formula compendium}
This appendix collects compact formulae that are useful during Q\&A.  They are intentionally redundant with the main text, because they provide a quick technical lookup while presenting.
\section{Kinematics and event reconstruction}
\begin{align}
s&=(p_{e^+}+p_{e^-})^2, & E_{\rm beam}&=\sqrt{s}/2,\\
p_{\rm miss}^\mu&=p_{e^+}^\mu+p_{e^-}^\mu-\sum_i p_i^\mu, & E_{\rm miss}&=p_{\rm miss}^0,\\
m_{ij}^2&=(p_i+p_j)^2, & m_{bjj}^2&=(p_b+p_{j_1}+p_{j_2})^2,\\
H_T&=\sum_i p_{T,i}, & \Delta R_{ij}&=\sqrt{(\eta_i-\eta_j)^2+(\phi_i-\phi_j)^2},\\
\cos\theta_{ij}&=\frac{\vec p_i\cdot\vec p_j}{|\vec p_i||\vec p_j|}, & \chi^2_W&=\frac{(m_{jj}-m_W)^2}{\sigma_W^2},\\
\chi^2_t&=\frac{(m_{bjj}-m_t)^2}{\sigma_t^2}, & \chi^2_{\rm tot}&=\chi^2_W+\chi^2_t.
\end{align}
\section{Electroweak couplings}
\begin{align}
Q_e&=-1, & T_3^{e_L}&=-\frac12, & T_3^{\nu_L}&=+\frac12,\\
g_L^f&=T_3^f-Q_f s_{\rm eff}^2, & g_R^f&=-Q_f s_{\rm eff}^2,\\
g_V^f&=g_L^f+g_R^f, & g_A^f&=g_L^f-g_R^f,\\
\Gamma_f&=N_c^f\frac{G_Fm_Z^3}{6\sqrt2\pi}\left[(g_V^f)^2+(g_A^f)^2\right](1+\delta_f),\\
A_f&=\frac{2g_V^fg_A^f}{(g_V^f)^2+(g_A^f)^2}, & A_{\rm FB}^{0,f}&=\frac34 A_eA_f,\\
R_b&=\Gamma_b/\Gamma_{\rm had}, & R_c&=\Gamma_c/\Gamma_{\rm had},\\
\Gamma_Z&=\sum_f\Gamma_f+\Gamma_{\rm inv}, & \Gamma_{\rm inv}&=N_\nu\Gamma_\nu^{\rm SM}.
\end{align}
\section{Higgs rates}
\begin{align}
\sigma_{ZH}&=\kappa_Z^2\sigma_{ZH}^{\rm SM}, & \Gamma_X&=\kappa_X^2\Gamma_X^{\rm SM},\\
\br_X&=\Gamma_X/\Gamma_H, & \mu_i^f&=\frac{\sigma_i\br_f}{(\sigma_i\br_f)_{\rm SM}},\\
\mu_i^f&=\frac{\kappa_i^2\kappa_f^2}{\kappa_H^2}, & \kappa_H^2&=\sum_j\kappa_j^2\br_j^{\rm SM},\\
\bar\kappa_H^2&=\frac{\kappa_H^2}{1-B_{\rm inv}-B_{\rm und}}, & \Gamma_H&=\bar\kappa_H^2\Gamma_H^{\rm SM},\\
\sigma(ZH)\br_{b\bar b}&\propto\frac{\kappa_Z^2\kappa_b^2}{\Gamma_H}, & \sigma(ZH)\br_{ZZ^*}&\propto\frac{\kappa_Z^4}{\Gamma_H},\\
\sigma(\nu\bar\nu H)\br_X&\propto\frac{\kappa_W^2\kappa_X^2}{\Gamma_H}, & m_{\rm recoil}^2&=s+m_Z^2-2\sqrt{s}E_Z.
\end{align}
\section{SMEFT expansion and fit structure}
\begin{align}
\mathcal L_{\rm SMEFT}&=\mathcal L_{\rm SM}+\sum_i\frac{C_i}{\Lambda^2}\mathcal O_i+\cdots,\\
\mathcal A&=\mathcal A_{\rm SM}+\sum_i\frac{C_i}{\Lambda^2}\mathcal A_i,\\
\sigma&=\sigma_{\rm SM}+\sum_i\frac{C_i}{\Lambda^2}\sigma_i^{\rm int}+\sum_{ij}\frac{C_iC_j}{\Lambda^4}\sigma_{ij}^{\rm quad},\\
\chi^2(\vec C)&=(\vec O_{\rm exp}-\vec O_{\rm th}(\vec C))^TV^{-1}(\vec O_{\rm exp}-\vec O_{\rm th}(\vec C)),\\
\delta_{\rm tot}^2&=\delta_{\rm stat}^2+\delta_{\rm syst}^2+\delta_{\rm th}^2+\delta_{\rm param}^2.
\end{align}
\section{Top FCNC partial widths and yields}
\begin{align}
N_S&=\sigma_S\mathcal L\epsilon_S, & N_B&=\sigma_B\mathcal L\epsilon_B,\\
Z_A&=\sqrt{2\left[(S+B)\ln\left(1+\frac{S}{B}\right)-S\right]}, & Z_{S/\sqrt B}&\simeq S/\sqrt B,\\
\Gamma(t\to q\gamma)&=\frac{\alpha}{4}m_t(|\lambda_L|^2+|\lambda_R|^2),\\
\Gamma(t\to qZ)_\kappa&=\frac{\alpha}{32s_W^2c_W^2}\frac{m_t^3}{m_Z^2}(1-r_Z)^2(2+r_Z)(|\kappa_L|^2+|\kappa_R|^2),\\
\br(t\to qX)&\simeq\Gamma(t\to qX)/\Gamma_t^{\rm SM}, & r_Z&=m_Z^2/m_t^2,\\
\mathcal L(\mu,\theta)&=\prod_{c,b}{\rm Pois}(n_{cb}|\mu s_{cb}(\theta)+b_{cb}(\theta))\prod_k\pi_k(\theta_k),\\
q_\mu&=-2\ln\frac{\mathcal L(\mu,\hat{\hat\theta}_\mu)}{\mathcal L(\hat\mu,\hat\theta)}, & CL_s&=\frac{p_{s+b}}{1-p_b}.
\end{align}

\part{Glossary and references}
\chapter{Glossary}
\begin{longtable}{p{0.18\textwidth}p{0.72\textwidth}}
\toprule
Term & Explanation \\
\midrule
FCC & Future Circular Collider. \\
FCC-ee & Electron--positron stage of the FCC integrated programme. \\
FCC-hh & Proton--proton stage of the FCC integrated programme. \\
HL-LHC & High-Luminosity Large Hadron Collider. \\
Tera-Z & FCC-ee Z-pole programme with about $10^{12}$ Z bosons. \\
EWPO & Electroweak precision observable. \\
SMEFT & Standard Model Effective Field Theory. \\
EFT & Effective Field Theory. \\
BSM & Beyond the Standard Model. \\
FCNC & Flavour-changing neutral current. \\
GIM & Glashow--Iliopoulos--Maiani suppression mechanism. \\
Higgsstrahlung & Associated production $\ee\to ZH$. \\
Recoil mass & Mass inferred from the system recoiling against the reconstructed Z boson. \\
VBF & Vector-boson fusion. \\
WW threshold & Energy region near $\ee\to W^+W^-$ threshold. \\
aTGC & Anomalous triple-gauge coupling. \\
IDEA & Innovative Detector for Electron-positron Accelerators. \\
EDM4HEP & Event Data Model for High-Energy Physics. \\
FCCSW & FCC software framework. \\
FCCAnalyses & FCC analysis framework for processing EDM4HEP samples. \\
UFO & Universal FeynRules Output model format. \\
CL$_s$ & Modified frequentist method for exclusion limits. \\
Asimov dataset & Representative dataset equal to the expectation, used for expected sensitivities. \\
\bottomrule
\end{longtable}

\begin{thebibliography}{99}
\bibitem{FCCFSR} FCC Collaboration, \emph{Future Circular Collider Feasibility Study Report Volume 1: Physics, Experiments, Detectors}, Eur. Phys. J. C 85 (2025) 1468, arXiv:2505.00272.
\bibitem{FCCCDR} FCC Collaboration, \emph{FCC Physics Opportunities: Future Circular Collider Conceptual Design Report Volume 1}, Eur. Phys. J. C 79 (2019) 474.
\bibitem{FCCManifesto} A. Blondel et al., \emph{The FCC integrated programme: a physics manifesto}, arXiv:2504.02634.
\bibitem{ECFAHET} ECFA Higgs, Electroweak and Top Factory Study, CERN Yellow Report Monograph 5/2025, arXiv:2506.15390.
\bibitem{PrecisionReach} T. Armadillo et al., \emph{New Physics Reach through Precision at Future Colliders: a Multi-Pronged Approach}, arXiv:2604.16596.
\bibitem{TopThreshold} M. M. Defranchis, J. de Blas, A. Mehta, M. Selvaggi and M. Vos, \emph{A detailed study on the prospects for a $t\bar t$ threshold scan in $e^+e^-$ collisions}, JHEP 11 (2025) 020, arXiv:2503.18713.
\bibitem{AguilarSaavedra2004} J. A. Aguilar-Saavedra, \emph{Top flavour-changing neutral interactions: Theoretical expectations and experimental detection}, Acta Phys. Polon. B35 (2004) 2695, hep-ph/0409342.
\bibitem{AguilarSaavedra2009} J. A. Aguilar-Saavedra, \emph{A minimal set of top anomalous couplings}, Nucl. Phys. B812 (2009) 181, arXiv:0811.3842.
\bibitem{FCCFCNCNote} P. Azzi, H. Khanpour, M. Mohammadi Najafabadi, E. Perez and S. Tizchang, \emph{Top quark FCNC anomalous couplings at the future collider FCC-ee}, FCC Physics and Experiments note, CERN repository record avxy7-mtz86.
\bibitem{FCCAfbNote} M. Cobal et al., \emph{Bottom quark forward-backward asymmetry at the future electron-positron collider FCC-ee}, FCC Physics and Experiments note, CERN repository record g97j3-q0w73.
\bibitem{IDEA} IDEA detector concept for FCC-ee, arXiv:2502.21223.
\bibitem{FCCAnalyses} FCCAnalyses documentation, \url{https://hep-fcc.github.io/FCCAnalyses/}.
\bibitem{Combine} CMS Higgs Combine Tool, \url{https://github.com/cms-analysis/HiggsAnalysis-CombinedLimit}.
\bibitem{ATLASProjection} ATLAS Collaboration, HL-LHC prospects for top FCNC, ATL-PHYS-PUB-2019-001.
\bibitem{CMSProjection} CMS Collaboration, HL-LHC prospects for top FCNC, CMS-CR-2018-094.
\end{thebibliography}
\end{document}
